% !PRINTED 466
\weberpage{466}

%% [running head, not body text -- commented out 2026-08-07] Page furniture (running
%% head, folio, signature) is excluded from this edition by standing policy.
%% \noindent 466 \hfill Neunzehnter Abschnitt. \hfill \S\,126.

Auf den Fall $m=1$ ist dies Verfahren nicht anwendbar, weil $\Phi_n(u,u)$ für $n=1$, $\omega=i$ (aber auch nur für diesen Wert) identisch verschwindet, für $m=1$ haben wir aber bereits früher gefunden (\S\,83):
\[
f(i)^{24}=64,
\]
so daß also auch in diesem Falle der ausgesprochene Satz gilt.

Wenn $m$ ungerade ist, so ist bei ungeradem $A$, $C$ der mittlere Koeffizient $B$ durch 4 teilbar; dagegen ist bei geradem $m$ unter der gleichen Voraussetzung $B$ nicht durch 4 teilbar. Durch die Substitution $\omega+1$ für $\omega$ erreichen wir aber auch im Fall eines geraden $m$, daß $B$ durch 4 teilbar wird; dann aber wird auch $C$ gerade. Durch die Vertauschung $(\omega,\omega+1)$ geht aber $f(\omega)^{24}$ in $-f_1(\omega)^{24}$ über, so daß wir also unserem Satz auch den folgenden Ausdruck geben können:

2. \gesperrt{Ist $\omega$ die Wurzel der quadratischen Gleichung:}
\[
A\omega^2+B\omega+C=0,
\]
\gesperrt{worin $A$ ungerade, $B$ gerade ist, so ist, je nachdem $D\equiv 4$ oder $D\equiv 0\pmod{8}$ ist, $f(\omega)^{24}$ oder $f_1(\omega)^{24}$ Klasseninvariante.}

Da $A$ und $C$ nicht beide gerade sein können, so sind nur drei Fälle möglich, von denen durch die Vertauschungen
\[
(\omega,\omega+1),\quad \left(\omega,-\frac{1}{\omega}\right)
\]
der zweite auf den ersten und der dritte auf den zweiten zurückgeführt wird. Wir haben also in diesen Fällen:
\weq{13}{
\begin{array}{lll}
A\equiv 1, & C\equiv 1\ (\mathrm{mod}\ 2), & \text{Klasseninvariante } f(\omega)^{24},\\
A\equiv 1, & C\equiv 0 & \text{\glqq}\quad\text{\glqq}\quad f_1(\omega)^{24},\\
A\equiv 0, & C\equiv 1 & \text{\glqq}\quad\text{\glqq}\quad f_2(\omega)^{24}.
\end{array}
}

Insbesondere ist also, wenn wir die Hauptform der Diskriminante $-4m$, $(1,0,m)$, zugrunde legen, $f(\sqrt{-m})^{24}$ bei ungeradem und $f_1(\sqrt{-m})^{24}$ bei geradem $m$ eine Klasseninvariante. Diese beiden Zahlen haben einen reellen positiven Wert, und sollen in den weiter unten folgenden Rechnungen vorzugsweise berücksichtigt werden.

Aus der Gleichung (1) schließen wir (Bd.\,II, \S\,154, 11.), da $j(\omega)$ eine ganze algebraische Zahl ist, daß auch $f(\omega)^{24}$ eine ganze algebraische Zahl sein muß, und die Form der Gleichung (1) zeigt, daß die Norm dieser Zahl eine Potenz von 2 ist.
